% Lab 2 — Constraining Methodology（原书 PDF p.19–32）
\bichapter{Lab 2：约束方法论}{Lab 2: Constraining Methodology}

\begin{objectiveBox}
完成本实验后，你应能够：
\begin{itemize}
  \item 通过检查约束完整性与未测试（untested）时序检查来验证约束；
  \item 在时序报告中识别并解读约束，包括：
  \begin{itemize}
    \item 时钟约束；
    \item 接口约束。
  \end{itemize}
\end{itemize}
\end{objectiveBox}

\begin{durationBox}
约 \textbf{30} 分钟。
\end{durationBox}

\section{概述}
\label{lab2:overview}

\noindent\textit{Overview}

本实验大致步骤如下：
\begin{enumerate}
  \item 验证约束；
  \item 在时序报告中解读约束。
\end{enumerate}

\subsection{相关文件与目录}
\label{lab2:files}

\noindent\textit{Relevant Files and Directories}

本实验全部文件位于主目录下的 \cmd{lab2_constraints} 目录中。

\begin{tabularx}{\textwidth}{@{}l X@{}}
\toprule
\textbf{路径 / 文件} & \textbf{说明} \\
\midrule
\cmd{lab2_constraints/} & 当前工作目录 \\
\cmd{orca_savesession/} & 用于恢复实验的会话目录 \\
\cmd{.synopsys_pt.setup} & 自动读取的 PT setup 文件 \\
\cmd{RUN.tcl} & 用于重建 \cmd{orca_savesession} 的运行脚本 \\
\bottomrule
\end{tabularx}

\subsection{答案 / 解答}
\label{lab2:answers-note}

\noindent\textit{Answers / Solutions}

本实验手册在每个 Lab 后半部分给出全部问题的答案与解答。
若答题需要帮助，或希望核对结果，请查阅本 Lab 后部的解答章节。

\section{操作说明}
\label{lab2:instructions}

\noindent\textit{Instructions}

\subsection{任务 1：验证约束}
\label{lab2:task1}

\noindent\textit{Task 1. Validate Constraints}

\begin{enumerate}
  \item 从实验 Unix 目录 \cmd{lab2_constraints} 启动 PrimeTime，并使用 \cmd{orca_savesession} 目录恢复 PrimeTime 会话。

\begin{noteBox}
\begin{itemize}
  \item 如需重建该会话，可使用：\\
        \cmd{pt_shell -f RUN.tcl | tee -i run.log}
  \item 执行 \cmd{RUN.tcl} 时出现的任何 \textbf{PARA-124} 错误，在本课程实验中可安全忽略。
\end{itemize}
\end{noteBox}

\begin{lstlisting}
unix% cd lab2_constraints
unix% pt_shell
pt_shell> restore_session orca_savesession
\end{lstlisting}

  \item 检查约束完整性：
\begin{lstlisting}
pt_shell> check_timing -verbose
\end{lstlisting}

\begin{questionBox}
\textbf{问题 1.}\ 设计中所有寄存器是否都已有时钟？\\
\textbf{问题 2.}\ 是否存在缺失的约束？你能解释原因吗？
\end{questionBox}

  \item 检查设计中未测试（untested）的时序检查：
\begin{lstlisting}
pt_shell> report_analysis_coverage
\end{lstlisting}

\begin{questionBox}
\textbf{问题 3.}\ 近三分之二的 setup/hold 检查处于 untested 状态！造成这一现象的两个原因是什么？\\
\textbf{问题 4.}\ 为何存在未行使（unexercised）的 \cmd{min_pulse_width} 检查？\\
\textbf{问题 5.}\ setup 与 hold 各有多少个 output delay 约束？这些约束是满足（met）还是违例（violated）？
\end{questionBox}
\end{enumerate}

\subsection{任务 2：分析输入延迟约束的时序报告}
\label{lab2:task2}

\noindent\textit{Task 2. Analyze a Timing Report For Input Delay Constraint}

\begin{enumerate}
  \item 生成施加于端口 \cmd{pad[0]} 的输入延迟约束报告：
\begin{lstlisting}
pt_shell> report_port -input_delay pad[0]
\end{lstlisting}

\begin{questionBox}
\textbf{问题 6.}\ 到达 \cmd{pad[0]} 的最小与最大到达时间（arrival time）分别是多少？\\
\textbf{问题 7.}\ 约束 \cmd{pad[0]} 的外部起点时钟名称是什么？
\end{questionBox}

  \item 生成从端口 \cmd{pad[0]} 开始的 setup 时序报告，并据此回答下列问题。
        可参考标注为 ``timing reports'' 的 job aid，回忆 \cmd{report_timing} 的合适开关。
\begin{lstlisting}
pt_shell> report_timing -from pad[0]
\end{lstlisting}

\begin{questionBox}
\textbf{问题 8.}\ 你使用了时序报告中的哪些行来确认该路径起始于端口 \cmd{pad[0]}，且为 setup 报告？\\
\textbf{问题 9.}\ 列出该时序报告中涉及的全部用户指定约束。\\
\textbf{问题 10.}\ 对于起点时钟 \cmd{PCI_CLK}，时钟延迟（clock latency）必须包含在何处？\\
\textbf{问题 11.}\ 描述端口 \cmd{pad[0]} 的方向（即它是 input、output 还是 inout 端口）。\\
\textbf{问题 12.}\ 描述该时序路径的终点（即它是 output 端口还是内部触发器）。
\end{questionBox}

  \item 从同一端口 \cmd{pad[0]} 生成另一份 setup 报告，同时显示计算得到的时钟网络延迟（clock network delay）详情。
        可参考 ``timing reports'' job aid，并善用历史命令。
\begin{lstlisting}
pt_shell> !rep -path full_clock
\end{lstlisting}

\begin{questionBox}
\textbf{问题 13.}\ 对于终点时钟 \cmd{PCI_CLK}，时钟源延迟（clock source latency）与时钟网络延迟各有多大？\\
\textbf{问题 14.}\ 时钟 \cmd{PCI_CLK} 在何处定义（时钟定义点）？
\end{questionBox}

  \item 生成从端口 \cmd{pad[0]} 开始的 hold 时序报告：
\begin{lstlisting}
pt_shell> report_timing -delay min -from pad[0]
\end{lstlisting}

\begin{questionBox}
\textbf{问题 15.}\ 输入外部延迟（input external delay）约束的值是否符合你的预期？
\end{questionBox}
\end{enumerate}

\subsection{任务 3：分析输出延迟约束的时序报告}
\label{lab2:task3}

\noindent\textit{Task 3. Analyze a Timing Report For Output Delay Constraint}

\begin{enumerate}
  \item 生成施加于端口 \cmd{pad[0]} 的输出延迟约束报告：
\begin{lstlisting}
pt_shell> report_port -output_delay pad[0]
\end{lstlisting}

\begin{questionBox}
\textbf{问题 16.}\ 该端口的最小与最大输出延迟约束分别是多少？\\
\textbf{问题 17.}\ 负的最小输出延迟约束将如何施加于该端口（即它会施加正向还是负向的 hold 要求）？\\
\textbf{问题 18.}\ 约束该端口的外部终点时钟名称是什么？
\end{questionBox}

  \item 生成一份以端口 \cmd{pad[0]} 为终点的 hold ``简短（short）'' 时序报告：
\begin{lstlisting}
pt_shell> report_timing -delay min -to pad[0] -path short
\end{lstlisting}

\begin{questionBox}
\textbf{问题 19.}\ 描述该时序路径的起点（即它是 input 端口还是内部触发器）。\\
\textbf{问题 20.}\ 该时序路径的 path group 是否符合你的预期？\\
\textbf{问题 21.}\ ``data required time'' 是否符合你的预期？
\end{questionBox}

  \item \textbf{［可选］} 施加下列约束，为 \cmd{pad[0]} 施加正向 hold 输出延迟约束，然后重新执行本任务中的步骤以观察影响：
\begin{lstlisting}
pt_shell> set_output_delay -min 1.0 -clock PCI_CLK pad[0]
\end{lstlisting}

  \item 退出 PrimeTime：
\begin{lstlisting}
pt_shell> quit
\end{lstlisting}
\end{enumerate}

\noindent 至此 Lab~2 完成。请返回课堂讲义。

\section{答案 / 解答}
\label{lab2:solutions}

\noindent\textit{Answers / Solutions}

\begin{answerBox}
\textbf{问题 1.}\ 是。在 \cmd{Information: Checking 'no_clock'.} 消息之后，没有报告任何时钟 pin。
\end{answerBox}

\begin{answerBox}
\textbf{问题 2.}\ 是。有 2 个 output 端口报告缺失 output delay：
\begin{lstlisting}
Warning: There are 2 endpoints which are not constrained for maximum
delay.

sd_CKn
sd_CK
\end{lstlisting}
上述警告可忽略——这 2 个是时钟输出端口（可用 \cmd{report_clock} 确认），不应施加 output delay 约束。
\end{answerBox}

\begin{answerBox}
\textbf{问题 3.}\ \textbf{constant\_disabled} 与 \textbf{false\_path}。
（可使用 \cmd{report\_analysis\_coverage -status untested -check \{setup hold\}} 确认。）
\end{answerBox}

\begin{answerBox}
\textbf{问题 4.}\ \cmd{min\_pulse\_width} 检查仅当 pin 上有时钟时才会被行使。
由于这些是诸如 set/clear 之类的非时钟异步 pin，未在其上定义时钟。
（可使用 \cmd{report\_analysis\_coverage -status untested -check min\_pulse\_width} 确认。）
\end{answerBox}

\begin{answerBox}
\textbf{问题 5.}\ 从 \cmd{report\_analysis\_coverage} 可见：setup 与 hold 各有 \textbf{75} 个 output delay 约束；
其中 setup 有 \textbf{9} 个违例，hold 有 \textbf{39} 个违例。
请确认所有 output 端口均已同时约束 setup 与 hold。
\begin{lstlisting}
pt_shell> restore_session orca_savesession
pt_shell> report_analysis_coverage
Type of Check    Total    Met         Violated    Untested
setup            9629     3575 ( 37%) 13  ( 0%)   6041 ( 63%)
hold             9629     3517 ( 37%) 71  ( 1%)   6041 ( 63%)
recovery         1316     1210 ( 92%) 0   ( 0%)   106  (  8%)
removal          1316     1204 ( 91%) 6   ( 0%)   106  (  8%)
min_period       20       20   (100%) 0   ( 0%)   0    (  0%)
min_pulse_width  7273     5957 ( 82%) 0   ( 0%)   1316 ( 18%)
out_setup        75       66   ( 88%) 9   ( 12%)  0    (  0%)
out_hold         75       36   ( 48%) 39  ( 52%)  0    (  0%)
All Checks       29333    15585( 53%) 138 (  0%)  13610( 46%)
\end{lstlisting}
\end{answerBox}

\begin{answerBox}
\textbf{问题 6.}\ 最小与最大到达时间分别为 \textbf{2\,ns} 与 \textbf{8\,ns}
（rise 与 fall 数据跳变在该端口使用相同约束值）。
\end{answerBox}

\begin{answerBox}
\textbf{问题 7.}\ 时钟名称为 \textbf{PCI\_CLK}。
\end{answerBox}

\begin{answerBox}
\textbf{问题 8.}
\begin{lstlisting}
pt_shell> report_timing -from pad[0]
Startpoint: pad[0] (input port clocked by PCI_CLK)
Endpoint:   I_ORCA_TOP/I_PCI_CORE/d_out_i_bus_reg[0]
            (rising edge-triggered flip-flop clocked by PCI_CLK)
Path Group: PCI_CLK
Path Type:  max
\end{lstlisting}
\end{answerBox}

\begin{answerBox}
\textbf{问题 9.}\ 时钟周期是一项约束；时钟 \cmd{PCI\_CLK} 为 propagated（非 ideal）；
输入外部延迟（来自 input delay 约束）也是用户指定约束。
\end{answerBox}

\begin{answerBox}
\textbf{问题 10.}\ 时钟网络延迟为零，因此表示外部时钟延迟的唯一位置是 input delay 约束中的输入外部延迟（input external delay）。
正确建模方式是：对 \cmd{set\_input\_delay} 使用 \opt{-network\_latency\_included} 与 \opt{-source\_latency\_included}。
\end{answerBox}

\begin{answerBox}
\textbf{问题 11.}\ 端口 \cmd{pad[0]} 是 \textbf{inout} 端口，因此它既是时序路径起点，也是时序路径终点。
\begin{lstlisting}
Point                                    Incr      Path
clock PCI_CLK (rise edge)                0.000     0.000
clock network delay (propagated)         0.000     0.000
input external delay                       8.000     8.000 r
pad[0] (inout)                             0.000     8.000 r
\end{lstlisting}
\end{answerBox}

\begin{answerBox}
\textbf{问题 12.}\ 终点是由 \cmd{PCI\_CLK} 时钟的上升沿触发 flip-flop
（实际上是一个具有 setup/hold 时序检查的时序模型，外观类似 flip-flop）。
\end{answerBox}

\begin{answerBox}
\textbf{问题 13.}\ 下方仅展示时序报告的 data required time 段。
源延迟（source latency）为 \textbf{0\,ns}；时钟网络延迟为 \textbf{0.768\,ns}（$15.768 - 15.000$）。
\begin{lstlisting}
pt_shell> !rep -path full_clock
clock PCI_CLK (rise edge)                15.000    15.000
clock source latency                       0.000    15.000
pclk (in)                                  0.000    15.000 r
...
I_ORCA_TOP/I_PCI_CORE/d_out_i_bus_reg[0]/CP (sdcrq1)
                                           0.088    15.768 r
library setup time                        -0.153    15.615
data required time                                  15.615
\end{lstlisting}
\end{answerBox}

\begin{answerBox}
\textbf{问题 14.}\ 时钟 \cmd{PCI\_CLK} 在输入端口 \cmd{pclk} 处定义。
时钟定义点将时钟源延迟与时钟网络延迟分开。
\end{answerBox}

\begin{answerBox}
\textbf{问题 15.}\ 是。由上文 \cmd{report\_port} 可知，输入外部延迟相对于 \cmd{PCI\_CLK} 上升沿应为 \textbf{2\,ns}，
时序报告确认了这一点：
\begin{lstlisting}
pt_shell> report_timing -delay min -from pad[0]
Startpoint: pad[0] (input port clocked by PCI_CLK)
Endpoint:   I_ORCA_TOP/I_PCI_CORE/d_out_i_bus_reg[0]
            (rising edge-triggered flip-flop clocked by PCI_CLK)
Path Group: PCI_CLK
Path Type:  min
...
clock PCI_CLK (rise edge)                0.000     0.000
clock network delay (propagated)         0.000     0.000
input external delay                       2.000     2.000 r
pad[0] (inout)                             0.000     2.000 r
\end{lstlisting}
\end{answerBox}

\begin{answerBox}
\textbf{问题 16.}
\begin{lstlisting}
pt_shell> report_port -output_delay pad[0]
Output Delay
Min    Max    Related    Related
Output Port  Rise  Fall  Rise  Fall  Clock    Pin
pad[0]      -1.00 -1.00  4.00  4.00  PCI_CLK  --
\end{lstlisting}
最小输出延迟约束为 \textbf{$-1.00$\,ns}，最大输出延迟约束为 \textbf{$4.00$\,ns}（rise/fall 相同）。
\end{answerBox}

\begin{answerBox}
\textbf{问题 17.}\ 课堂讲义指出：负的 hold output delay 约束会施加\textbf{正向} hold 要求。
\end{answerBox}

\begin{answerBox}
\textbf{问题 18.}\ 端口 \cmd{pad[0]} 相对于 \textbf{PCI\_CLK} 施加约束。
\end{answerBox}

\begin{answerBox}
\textbf{问题 19.}\ 时序路径起点是内部 flip-flop：
\begin{lstlisting}
pt_shell> report_timing -delay min -to pad[0]
Startpoint: I_ORCA_TOP/I_PCI_CORE/pad_out_buf_reg[0]
            (rising edge-triggered flip-flop clocked by PCI_CLK)
Endpoint:   pad[0] (output port clocked by PCI_CLK)
\end{lstlisting}
\end{answerBox}

\begin{answerBox}
\textbf{问题 20.}\ 是。path group 为 \textbf{PCI\_CLK}，与外部捕获时钟名称一致。
\end{answerBox}

\begin{answerBox}
\textbf{问题 21.}\ 是。捕获时钟边为 0\,ns，这对 hold 是合适的。
hold 要求为正向 \textbf{1\,ns}，propagated 时钟网络延迟为 0\,ns：
\begin{lstlisting}
clock PCI_CLK (rise edge)                0.000     0.000
clock network delay (propagated)         0.000     0.000
output external delay                      1.000     1.000
data required time                                   1.000
\end{lstlisting}
\end{answerBox}

\noindent\textbf{［可选步骤］} 施加 \cmd{set\_output\_delay -min 1.0 -clock PCI\_CLK pad[0]} 后：
\begin{answerBox}
\begin{lstlisting}
pt_shell> report_port -output_delay pad[0]
Output Delay
Min    Max    Related    Related
Output Port  Rise  Fall  Rise  Fall  Clock    Pin
pad[0]       1.00  1.00  4.00  4.00  PCI_CLK  --

pt_shell> report_timing -to pad[0] -delay min -path short
...
clock network delay (propagated)         0.772     0.772
...
data arrival time                                  3.445
clock PCI_CLK (rise edge)                0.000     0.000
clock network delay (propagated)         0.000     0.000
output external delay                     -1.000    -1.000
data required time                                  -1.000
slack (MET)                                          4.445
\end{lstlisting}
由上可见：对 output delay 使用正向 hold 约束会\textbf{增大}正向 slack。
这证实：对 output delay 指定负 hold 约束，实际上是在 output 端口上指定 hold 要求。
\end{answerBox}

